O cronograma original precisou ser ajustado para a realização do software. A adoção de um desenvolvimento orientado por Inteligência Artificial (IA) exigiu um período adicional dedicado ao estudo e à compreensão da metodologia, que não havia sido previsto no planejamento inicial. Esse investimento representou uma troca deliberada, o tempo dedicado ao aprendizado das técnicas de IA foi compensado pela redução do tempo de codificação nas etapas seguintes, tornando o desenvolvimento mais ágil. A Tabela~\ref{tab:cronograma} apresenta o cronograma atualizado, com os épicos do projeto e suas respectivas datas de entrega.

\renewcommand{\arraystretch}{1.5}
\begin{table}[H]
    \centering
    \caption{Cronograma de atividades (TCC 2).}
    \label{tab:cronograma}

    \begin{tabularx}{\textwidth}{|X|c|c|c|c|}
        \hline
        \textbf{Atividade} & \textbf{Set} & \textbf{Out} & \textbf{Nov} & \textbf{Dez} \\
        \hline
        Revisão da proposta de solução & X &   &   &   \\
        \hline
        Desenvolvimento da aplicação   & X & X & X &   \\
        \hline
        Teste do aplicativo            &   & X  & X &  \\
        \hline
        Geração de versão do aplicativo &  &   & X &   \\
        \hline
        Realização de correções        &   &   & X &  \\
        \hline
        Avaliação de resultados        &   &   & X  & X \\
        \hline
        Apresentação do TCC 2          &   &   &   & X  \\
        \hline
    \end{tabularx}

    \vspace{0.5em}
    \small Fonte: autoria própria.
\end{table}